\section{Related Work}
\label{sec:related}

\subsection{Vision-Language-Action Models}
VLA models extend pretrained vision-language models with action prediction heads, enabling robots to follow natural-language instructions from visual observations. Early work such as Gato~\cite{reed2022gato} and RT-1~\cite{brohan2023rt1} demonstrated that transformer-based policies can generalize across many tasks from diverse robot datasets. The Open X-Embodiment project~\cite{padalkar2023openx} aggregated large-scale cross-embodiment data, enabling generalist models such as Octo~\cite{ghosh2024octo}. OpenVLA~\cite{kim2024openvla} demonstrated that a 7\,B-parameter LLM backbone could be fine-tuned for generalizable manipulation across diverse tasks. OpenVLA-OFT~\cite{kim2025openvlaoft} introduced Optimized Fine-Tuning (OFT) with parallel decoding and an action chunking head, establishing strong LIBERO baselines. $\pi_0$~\cite{black2024pi0} adopts a flow-matching formulation, while SmolVLA~\cite{shukor2025smolvla} and CoT-VLA~\cite{zhao2025cotvla} explore efficient and reasoning-augmented VLA designs respectively. UniVLA~\cite{bu2025univla} and GR00T N1~\cite{nvidia2025groot} further advance cross-embodiment generalization via latent action spaces and humanoid-oriented pretraining. On the efficiency axis, BitVLA~\cite{wang2025bitvla} builds on the BitNet~\cite{wang2023bitnet} and BitNet b1.58~\cite{ma2024bitnet158} quantization framework to compress model weights to 1.58-bit precision. Using a BitSigLIP~\cite{zhai2023siglip} vision encoder, it achieves 96.0\% average success on LIBERO with only 3\,B parameters and 1.4\,GB memory---competitive with models several times larger.

TAR is orthogonal to all of the above architectural choices. It modifies only the training objective and is applicable to any VLA model that predicts action chunks, regardless of backbone size or quantization scheme.

\subsection{Action Chunking and Temporal Consistency}
Action chunking was popularized by ACT~\cite{zhao2023act}, which showed that predicting $T$-step sequences reduces compounding errors and improves dexterity on fine-grained bimanual tasks. Diffusion Policy~\cite{chi2023diffusion} generates action trajectories over a receding horizon, with temporal coherence emerging implicitly from the denoising process. OpenVLA-OFT~\cite{kim2025openvlaoft} adopts a parallel action head that predicts full chunks in a single forward pass.

To mitigate discontinuities at chunk boundaries, ACT~\cite{zhao2023act} proposes temporal ensembling: overlapping windows are executed and averaged online. While effective, ensembling increases inference latency and complicates deployment. TAR takes a complementary approach---eliminating boundary jumps at training time so that a single, non-overlapping chunk execution suffices.

\subsection{Regularization in Robot Learning}
Behavior cloning (BC) and its variants~\cite{ross2011reduction} form the backbone of imitation learning for robot manipulation~\cite{mandlekar2021matters}, yet they treat all timesteps uniformly without regard to temporal structure. Regularization in imitation learning typically operates on model parameters (weight decay, dropout) or the data distribution (augmentation, noise injection)~\cite{kim2025openvlaoft,wang2025bitvla}. BC-Z~\cite{jang2022bcz} and RT-2~\cite{brohan2023rt2} train with L1/L2 regression over action sequences, which implicitly encourages intra-chunk smoothness but places no explicit constraint on boundary steps. Closest to our work,~\cite{marino2024gradient} proposes gradient-based regularization to penalize abrupt action changes in reinforcement learning. Our approach differs in targeting the specific structural discontinuity at chunk boundaries in imitation-learning-based VLA models---a problem that does not arise in single-step RL policies.

TAR is the first to identify the training-time gap between intra-chunk and inter-chunk smoothness and to close it with a targeted boundary-velocity regularizer. Unlike post-hoc filtering or inference-time ensembling, TAR shapes the learned policy distribution itself, yielding smoother trajectories with zero inference overhead.
